\answer{ex:16:1}{La sangre es un circuito RC con $\tau=t_{1/2}/\ln2=14.4$~min; máximo sobre mínimo, $2^{T/t_{1/2}}=64$; el armónico fundamental se atenúa a $0.55$; la razón es $2$ con $t_{1/2}=T=60$~min.}
\answer{ex:16:2}{$\varphi^*=\arcsin\bigl((\omega-\omega_0)/K_\varphi\bigr)\approx41$~min ($\omega-\omega_0\approx0.047$~rad/día); tiempo de relajación $1/(K_\varphi\cos\varphi^*)\approx3.9$~días; tras un desplazamiento de $+6$ y $-6$~h, $7.7$ y $7.9$~días.}
\answer{ex:16:3}{$K_c(n)=\sec^n(\pi/n)$: $8$ y $2.37$, error del $11\%$ y del $30\%$; cuando $n\to\infty$, $K_c\to1$ y el error $\to50\%$.}
\answer{ex:16:4}{Las órdenes extremas (acelerador $80$, freno $0$) dan $f=120$ en $0.76$~s; “solo acelerador”, $2.33$~s, el triple.}
\answer{ex:16:5}{$3\arctan(\omega\tau)+\omega d=\pi$, $K_c=(1+\omega^2\tau^2)^{3/2}$: $K_c=8$, $3.54$, $2.50$, $1.76$ y periodo $3.63$, $5.46$, $6.86$, $9.27\,\tau$; con $0.9K_c$ las oscilaciones se amortiguan, con $1.1K_c$ se sostienen.}
\answer{ex:16:6}{$0.60\bar c$, $0.87\bar c$, $0.90\bar c$, límite $0.91\bar c$; con concentración constante la respuesta es nula: un receptor así solo lee porciones.}
\answer{ex:16:7}{Ejercicio abierto. Con realimentación hay ritmo solo si $K>8$; el periodo es proporcional a la $\tau$ de las etapas y casi no depende de $K$ ($3.64\,\tau$ con $K=8.5$, $4.23\,\tau$ con $K=40$): variar $K$ en un factor de $1.5$ desplaza el periodo un $2$--$4\%$, menos que el error de medida. Lo decide abrir el lazo (una hormona final constante a la entrada de la primera etapa suprime ese ritmo) o la respuesta a un breve aporte de hormona en distintas fases del ciclo.}
